\documentclass[11pt]{article}

\usepackage[margin=1in]{geometry}
\usepackage{amsmath,amssymb,amsfonts,bm,mathtools}
\usepackage{booktabs}
\usepackage{enumitem}
\usepackage[hidelinks]{hyperref}

\title{Theoretical Derivations for a Lightweight Axisymmetric Taylor-Cone Electrospray Solver}
\author{ISEF Physics Project --- Theory Deliverable}
\date{\today}

\newcommand{\eps}{\varepsilon}
\newcommand{\vE}{\mathbf{E}}
\newcommand{\vD}{\mathbf{D}}
\newcommand{\vJ}{\mathbf{J}}
\newcommand{\vU}{\mathbf{u}}
\newcommand{\vn}{\mathbf{n}}
\newcommand{\vT}{\mathbf{T}}
\newcommand{\grad}{\nabla}
\newcommand{\divg}{\nabla\!\cdot}
\newcommand{\dd}{\,\mathrm{d}}
\newcommand{\abs}[1]{\left|#1\right|}
\newcommand{\BoE}{\mathrm{Bo}_E}

\begin{document}
\maketitle

\begin{abstract}
This document contains the theoretical derivations needed before constructing the computational model for an axisymmetric Taylor-cone / electrospray-onset solver.  The goal is to define the continuous mathematical problem with enough precision that a separate implementation document can translate it into code.  We derive the classical Taylor cone angle, identify the apex singularity, formulate axisymmetric Laplace and Poisson electrostatics, derive the Maxwell--Young--Laplace interfacial stress balance, define curvature in axisymmetric geometry, introduce effective space-charge shielding closures, and state the verification targets required before physical validation.
\end{abstract}

\tableofcontents

\section{Purpose of the theoretical model}

The product we want is a lightweight solver that predicts Taylor-cone geometry and onset-like voltage trends for an electrospray emitter.  The solver should be mathematically anchored in the known Taylor-cone limit and then extend that limit to include effective space-charge shielding near the cone apex.

The theory phase must answer the following questions before coding:
\begin{enumerate}[label=\arabic*.]
    \item What is the ideal analytical solution we must reproduce?
    \item What equations govern the electric potential in axisymmetric coordinates?
    \item How does the electric field produce stress on the liquid surface?
    \item How does surface tension enter through curvature?
    \item How do we define a residual that measures whether a candidate interface is physically balanced?
    \item How can space charge be introduced without claiming a full plasma/electrospray emission model?
    \item What checks prove that the model is theoretically consistent before experiments begin?
\end{enumerate}

The model is deliberately reduced-order.  The first working product is not a full industrial multiphase CFD solver.  It is an open, sparse, axisymmetric electrostatic--capillary solver with optional Poisson space-charge shielding.

\section{Model hierarchy}

We use a staged hierarchy:

\begin{center}
\begin{tabular}{@{}clp{0.55\textwidth}@{}}
\toprule
Level & Model & Role \\
\midrule
0 & Classical Taylor cone & Analytical benchmark: recover the \(49.3^\circ\) cone half-angle. \\
1 & Axisymmetric Laplace electrostatics & Solve charge-free potential and field. \\
2 & Maxwell--Young--Laplace residual & Evaluate capillary/electric normal stress balance on an interface. \\
3 & Free-boundary equilibrium & Adjust an interface shape until the residual is minimized. \\
4 & Poisson space charge & Study how charge density shields the apex field. \\
5 & Leaky dielectric / transport & Optional later extension, not required for the first solver product. \\
\bottomrule
\end{tabular}
\end{center}

The first four levels define the core theoretical product.

\section{Geometry and coordinates}

We assume axisymmetry and use cylindrical coordinates
\begin{equation}
    (r,\theta,z), \qquad r\ge 0,
\end{equation}
with no \(\theta\)-dependence.  The gas/air domain is \(\Omega_g\), the liquid domain is \(\Omega_l\), and the liquid--gas interface is \(\Gamma\).  The electric potential is \(\phi\), and
\begin{equation}
    \vE=-\grad\phi.
\end{equation}

An interface can be represented as a graph
\begin{equation}
    r=R(z),
\end{equation}
or parametrically as
\begin{equation}
    s\mapsto (r(s),z(s)).
\end{equation}
For the first model, graph form is sufficient except near possible overturning, which is not part of the initial onset calculation.

\section{Classical Taylor cone derivation}

\subsection{Laplace equation outside a conducting cone}

For an ideal perfectly conducting liquid in a charge-free gas, the gas-domain potential satisfies
\begin{equation}
    \nabla^2\phi=0.
\end{equation}
Near the cone apex, spherical coordinates \((\rho,\vartheta)\) are natural, where \(\vartheta\) is measured from the symmetry axis.  The axisymmetric Laplace operator is
\begin{equation}
\nabla^2\phi
=
\frac{1}{\rho^2}\frac{\partial}{\partial\rho}
\left(\rho^2\frac{\partial\phi}{\partial\rho}\right)
+
\frac{1}{\rho^2\sin\vartheta}\frac{\partial}{\partial\vartheta}
\left(\sin\vartheta\frac{\partial\phi}{\partial\vartheta}\right).
\end{equation}
A separated axisymmetric harmonic solution is
\begin{equation}
    \phi(\rho,\vartheta)=A\rho^\nu P_\nu(\cos\vartheta),
    \label{eq:separated}
\end{equation}
where \(P_\nu\) is a Legendre function of degree \(\nu\), not necessarily an integer.

\subsection{Stress-scaling argument}

The electric field scales as
\begin{equation}
    E\sim \abs{\grad\phi}\sim \rho^{\nu-1}.
\end{equation}
The electrostatic pressure scales as
\begin{equation}
    p_E\sim \eps_0E^2\sim \rho^{2\nu-2}.
\end{equation}
For a cone, the curvature scale is \(1/\rho\), so capillary pressure scales as
\begin{equation}
    p_\gamma\sim \frac{\gamma}{\rho}.
\end{equation}
Balancing electric pressure and capillary pressure requires
\begin{equation}
    2\nu-2=-1,
\end{equation}
after which
\begin{equation}
    \nu=\frac12.
\end{equation}
Thus the local Taylor potential is
\begin{equation}
    \phi(\rho,\vartheta)=A\rho^{1/2}P_{1/2}(\cos\vartheta).
\end{equation}

\subsection{Cone angle condition}

A perfect conductor is an equipotential.  If the cone surface is \(\vartheta=\vartheta_0\), the angular part must vanish:
\begin{equation}
    P_{1/2}(\cos\vartheta_0)=0.
    \label{eq:taylor_root}
\end{equation}
The relevant root is approximately
\begin{equation}
    \vartheta_0\approx130.7^\circ.
\end{equation}
The physical cone semi-vertical angle is therefore
\begin{equation}
    \alpha_T=180^\circ-\vartheta_0\approx49.3^\circ.
\end{equation}
This is the most important analytical benchmark for the solver.

\subsection{Apex singularity}

With \(\nu=1/2\),
\begin{equation}
    E\sim\rho^{-1/2}.
\end{equation}
Therefore
\begin{equation}
    E\to\infty \quad \text{as}\quad \rho\to0.
\end{equation}
This singularity is unphysical.  In reality it is regularized by finite apex radius, charge emission, space charge, finite conductivity, viscosity, and molecular-scale effects.  The computational model will regularize it through finite grid/interface resolution and through Poisson space-charge shielding.

\section{Axisymmetric electrostatics}

For an axisymmetric potential \(\phi(r,z)\),
\begin{equation}
    \nabla^2\phi
    =
    \frac{1}{r}\frac{\partial}{\partial r}
    \left(r\frac{\partial\phi}{\partial r}\right)
    +
    \frac{\partial^2\phi}{\partial z^2}.
    \label{eq:axi_operator}
\end{equation}

\subsection{Laplace model}

The charge-free gas model is
\begin{equation}
    \frac{1}{r}\frac{\partial}{\partial r}
    \left(r\frac{\partial\phi}{\partial r}\right)
    +\frac{\partial^2\phi}{\partial z^2}=0.
    \label{eq:laplace_axi}
\end{equation}

\subsection{Poisson model}

With gas-domain charge density \(\rho_e\),
\begin{equation}
    \frac{1}{r}\frac{\partial}{\partial r}
    \left(r\frac{\partial\phi}{\partial r}\right)
    +\frac{\partial^2\phi}{\partial z^2}
    =-\frac{\rho_e(r,z)}{\eps_0}.
    \label{eq:poisson_axi}
\end{equation}

\subsection{Boundary conditions}

The essential boundary conditions are:
\begin{align}
    \phi&=V_0 &&\text{on powered conductor or ideal liquid surface},\\
    \phi&=0 &&\text{on grounded extractor electrode},\\
    \partial_r\phi&=0 &&\text{on the symmetry axis } r=0,\\
    \partial_n\phi&=0\text{ or }\phi=\phi_\infty &&\text{on artificial far boundaries}.
\end{align}

The axis condition is required for regularity.  The far-boundary condition is a numerical approximation and must be tested by moving the boundary farther away.

\section{Maxwell stress and electric pressure}

The electrostatic Maxwell stress tensor is
\begin{equation}
    \vT_E=\eps\left(\vE\otimes\vE-\frac12E^2\mathbf I\right).
\end{equation}
For a perfect conductor at electrostatic equilibrium, the electric field inside the liquid is zero and the tangential gas-side electric field vanishes on the surface.  The normal electric pressure on the liquid surface reduces to
\begin{equation}
    p_E=\frac{\eps_g}{2}E_n^2,
    \qquad
    E_n=\vE_g\cdot\vn.
\end{equation}

For leaky dielectric extensions, the complete traction jump would use both gas and liquid Maxwell stresses:
\begin{equation}
    \vn\cdot(\vT_{E,g}-\vT_{E,l})\vn.
\end{equation}
The first product uses the conducting limit as the basic model.

\section{Capillarity and curvature}

Surface tension produces a pressure jump proportional to mean curvature:
\begin{equation}
    \Delta p_\gamma=\gamma\kappa.
\end{equation}
For an axisymmetric graph \(r=R(z)\), using the convention that the gas is toward increasing \(r\), one convenient normal is
\begin{equation}
    \vn=\frac{(1,-R_z)}{\sqrt{1+R_z^2}}.
\end{equation}
The axisymmetric curvature is
\begin{equation}
    \kappa
    =
    \frac{1}{R\sqrt{1+R_z^2}}
    -
    \frac{R_{zz}}{(1+R_z^2)^{3/2}}.
    \label{eq:curvature}
\end{equation}
The first term is azimuthal curvature and the second is meridional curvature.

\section{Young--Laplace--Maxwell balance}

The quasi-static conducting-liquid normal stress balance is
\begin{equation}
    \gamma\kappa=\Delta p+\frac{\eps_g}{2}E_n^2,
    \label{eq:ylm}
\end{equation}
where \(\Delta p\) is a constant pressure offset.  Because curvature sign conventions vary, the implementation must choose one convention and maintain it consistently.

A useful residual is
\begin{equation}
    \mathcal R(s)
    =
    \gamma\kappa(s)-\Delta p-\frac{\eps_g}{2}E_n(s)^2.
    \label{eq:residual}
\end{equation}
A candidate equilibrium interface should satisfy
\begin{equation}
    \mathcal R(s)\approx0
\end{equation}
along \(\Gamma\).  The RMS residual is
\begin{equation}
    \mathcal R_{\rm RMS}
    =
    \left(\frac{1}{L_\Gamma}\int_\Gamma \mathcal R^2\dd s\right)^{1/2}.
\end{equation}
This residual is the bridge between theory and computation: it tells the code how physically balanced a guessed interface is.

\section{Space-charge shielding theory}

Poisson's equation requires a closure for \(\rho_e\).  The project should use a hierarchy of closures rather than claiming a complete microscopic emission model.

\subsection{Closure A: prescribed localized charge cloud}

The simplest shielding model is
\begin{equation}
    \rho_e(r,z)
    =
    \rho_0\exp\left[-\frac{(r-r_a)^2+(z-z_a)^2}{2\ell^2}\right],
    \label{eq:gaussian}
\end{equation}
where \((r_a,z_a)\) is the apex location and \(\ell\) is a shielding length.  This model tests how a localized charge distribution changes the electric field.

\subsection{Closure B: threshold-activated effective charge}

A more physical effective closure activates charge only above a field scale \(E_c\):
\begin{equation}
    \rho_e
    =
    \rho_{\max}
    \left[1-\exp\left(-\frac{\abs{\vE}-E_c}{E_s}\right)\right]_+,
    \label{eq:threshold}
\end{equation}
where \([x]_+=\max(x,0)\), \(E_s\) smooths the transition, and \(\rho_{\max}\) is a saturation density.  Since \(\rho_e\) depends on \(\vE=-\grad\phi\), this gives a nonlinear Poisson problem.

A fixed-point iteration is
\begin{align}
    \nabla^2\phi^{(k+1)}&=-\rho_e^{(k)}/\eps_0,\\
    \vE^{(k+1)}&=-\grad\phi^{(k+1)},\\
    \widetilde\rho_e^{(k+1)}&=F(\abs{\vE^{(k+1)}}),\\
    \rho_e^{(k+1)}&=(1-\omega)\rho_e^{(k)}+\omega\widetilde\rho_e^{(k+1)},
\end{align}
where \(0<\omega\le1\) is an under-relaxation parameter.

\subsection{Later extension: drift-diffusion-Poisson}

A more complete gas charge model solves
\begin{equation}
    \frac{\partial\rho_e}{\partial t}+\divg\vJ=S,
    \qquad
    \vJ=\mu_c\rho_e\vE-D_c\grad\rho_e,
\end{equation}
coupled with Poisson's equation.  This is a later extension because it introduces mobility, diffusion, and source parameters that may not be known in the first stage.

\section{Leaky-dielectric extension}

A real liquid may not be a perfect conductor.  In a leaky-dielectric model,
\begin{align}
    \divg(\eps_g\grad\phi_g)&=-\rho_{e,g} &&\text{in }\Omega_g,\\
    \divg(\eps_l\grad\phi_l)&=-\rho_{e,l} &&\text{in }\Omega_l.
\end{align}
If Ohmic conduction dominates inside the liquid,
\begin{equation}
    \vJ_l=\sigma_l\vE_l=-\sigma_l\grad\phi_l,
\end{equation}
and charge conservation gives
\begin{equation}
    \divg(\sigma_l\grad\phi_l)=0.
\end{equation}
At the interface,
\begin{align}
    \phi_g&=\phi_l,\\
    \vn\cdot(\eps_g\vE_g-\eps_l\vE_l)&=q_s.
\end{align}
Surface charge conservation has the general form
\begin{equation}
    \frac{\partial q_s}{\partial t}
    +\nabla_s\cdot(q_s\vU_s)
    +q_s(\nabla_s\cdot\vU_s)
    =
    \vn\cdot(\sigma_l\vE_l-\sigma_g\vE_g)-j_{\rm emit}.
\end{equation}
This is included to show the continuum target, but the first working solver should not rely on solving this full system.

\section{Nondimensionalization}

Let \(L\) be a characteristic length scale, \(V_0\) the voltage scale, and \(\gamma\) the surface tension.  Define
\begin{equation}
    \hat r=\frac{r}{L},\qquad
    \hat z=\frac{z}{L},\qquad
    \hat\phi=\frac{\phi}{V_0}.
\end{equation}
Then
\begin{equation}
    \vE=-\frac{V_0}{L}\hat\grad\hat\phi.
\end{equation}
The dimensionless charge density is
\begin{equation}
    \hat\rho_e=\frac{L^2\rho_e}{\eps_0V_0}.
\end{equation}
The electric-to-capillary stress ratio is
\begin{equation}
    \BoE=\frac{\eps_0V_0^2}{\gamma L}.
\end{equation}
This parameter controls how strongly electric stress competes with surface tension.

\section{Theoretical deliverables before coding is considered successful}

Before the solver is treated as a working theoretical product, it must demonstrate:
\begin{enumerate}[label=\arabic*.]
    \item \textbf{Taylor limit:} recovery of \(\alpha_T\approx49.3^\circ\) in the zero-space-charge conducting limit.
    \item \textbf{Correct operator:} implementation of the axisymmetric Laplace/Poisson operator, including axis regularity.
    \item \textbf{Stress consistency:} Maxwell pressure scales as \(V_0^2\).
    \item \textbf{Curvature consistency:} curvature signs and magnitudes are stable for known curves.
    \item \textbf{Residual meaning:} equilibrium-like shapes have smaller \(\mathcal R_{\rm RMS}\) than deliberately poor shapes.
    \item \textbf{Space-charge effect:} the Poisson model changes peak field and can represent shielding relative to the Laplace case.
    \item \textbf{Parameter transparency:} every non-ideal parameter, such as \(E_c\), \(\rho_{\max}\), or \(\ell\), is declared as a model parameter rather than hidden.
\end{enumerate}

\section{Conclusion}

The theory required for the working product is now defined.  The solver should begin from the classical Taylor cone, solve axisymmetric Laplace/Poisson electrostatics, compute electric fields, evaluate Maxwell pressure and capillary curvature on an interface, and use the Young--Laplace--Maxwell residual as the central physical diagnostic.  Space-charge shielding should be introduced as an explicit Poisson closure, first phenomenologically and then through threshold-activated nonlinear iteration.  This gives a rigorous theoretical foundation for the separate coding implementation outline.

\begin{thebibliography}{9}
\bibitem{Taylor1964}
G. I. Taylor, ``Disintegration of water drops in an electric field,'' \emph{Proceedings of the Royal Society of London. Series A}, vol. 280, no. 1382, pp. 383--397, 1964.

\bibitem{MelcherTaylor1969}
J. R. Melcher and G. I. Taylor, ``Electrohydrodynamics: A review of the role of interfacial shear stresses,'' \emph{Annual Review of Fluid Mechanics}, vol. 1, pp. 111--146, 1969.

\bibitem{Saville1997}
D. A. Saville, ``Electrohydrodynamics: The Taylor--Melcher leaky dielectric model,'' \emph{Annual Review of Fluid Mechanics}, vol. 29, pp. 27--64, 1997.
\end{thebibliography}

\end{document}
